% !TEX root = ../main.tex

\ifprmresultsready

%----------------------------------------------------------------------
\subsection{Effects of defect-aware model components}
\label{sec:factorial_results}
%----------------------------------------------------------------------

Validation MAE selected the all-enabled \PRMSelectedVariant{} variant.  Paired
contrasts resolved lower errors for gated readout and the composite
pre-normalized, distance-gated local block, whereas the independent main effect
of environment enrichment was inconclusive.  The complete factorial is
reported in Table~\ref{tab:factorial}, with its paired-effect visualization in
Supplemental Fig.~S2~\cite{supplemental}.  The selected
variant had mean validation MAE
\PRMSelectedValidationMAE\eV{} (95\% bootstrap interval
[\PRMSelectedValidationMAELow, \PRMSelectedValidationMAEHigh]\eV{}) across
the five paired repeats and test MAE
\PRMSelectedTestMAE\eV{} [\PRMSelectedTestMAELow,
\PRMSelectedTestMAEHigh]\eV{}.

% Auto-generated; do not edit.
\begin{table*}[t]
\caption{Paired $2^3$ architecture factorial. Values are mean MAE with
95\% bootstrap intervals over five paired repeats, in eV. G denotes gated
readout, E defect-environment enrichment, and P the composite pre-norm
distance-gated local block. Bold marks the validation-selected variant; test
MAE was not used in its selection.}
\label{tab:factorial}
\centering
\begin{tabular}{lccc cc}
\toprule
Variant & G & E & P & Validation MAE & Locked test MAE \\
\midrule
\texttt{g000} & 0 & 0 & 0 & 0.422 [0.413, 0.430] & 0.419 [0.410, 0.433] \\
\texttt{g001} & 0 & 0 & 1 & 0.394 [0.375, 0.413] & 0.402 [0.393, 0.413] \\
\texttt{g010} & 0 & 1 & 0 & 0.404 [0.390, 0.417] & 0.412 [0.405, 0.421] \\
\texttt{g011} & 0 & 1 & 1 & 0.401 [0.396, 0.407] & 0.410 [0.394, 0.429] \\
\texttt{g100} & 1 & 0 & 0 & 0.401 [0.389, 0.413] & 0.413 [0.407, 0.420] \\
\texttt{g101} & 1 & 0 & 1 & 0.374 [0.357, 0.388] & 0.387 [0.375, 0.400] \\
\texttt{g110} & 1 & 1 & 0 & 0.404 [0.396, 0.410] & 0.405 [0.394, 0.415] \\
\textbf{\texttt{g111}} & 1 & 1 & 1 & \textbf{0.371 [0.357, 0.385]} & \textbf{0.382 [0.377, 0.387]} \\
\bottomrule
\end{tabular}
\end{table*}


The orthogonal contrasts separate evidence for the three modules
from variation between random partitions.  Because P jointly changes
normalization placement and distance gating, its contrast supports or rejects
that composite block only; it does not identify either subchange as the causal
source of an effect.
% Auto-generated from canonical factorial claims; do not edit.
On validation MAE, the enabled-minus-disabled main-effect contrasts were $G$: $\Delta=-0.018$ [-0.020, -0.016]~eV (supported reduction); $E$: $\Delta=-0.003$ [-0.010, +0.004]~eV (inconclusive); $P$: $\Delta=-0.022$ [-0.029, -0.018]~eV (supported reduction). Directional validation intervals were also observed for $G:P$: $\Delta=-0.007$ [-0.011, -0.003]~eV (supported reduction); $E:P$: $\Delta=+0.005$ [+0.000, +0.010]~eV (supported increase). All other prespecified interaction intervals included zero. The repeat-level contrast agreed with the pooled point-estimate direction in $G$ 5/5, $P$ 5/5, $G:P$ 5/5, $E:P$ 4/5 paired repeats. With five paired repeats, the percentile intervals and sign counts are descriptive uncertainty summaries rather than large-sample hypothesis tests.


%----------------------------------------------------------------------
\subsection{Interpolation, chemical transfer, and low-energy behavior}
\label{sec:transfer_results}
%----------------------------------------------------------------------

Chemical transfer was strongly axis dependent.  Mean absolute error was
\PRMPairCvDARTMAE\eV{} when complete host--impurity pairs were withheld, versus
\PRMIdCvDARTMAE\eV{} under random cross-validation.  Withholding complete hosts
or impurities increased MAE to \PRMHostCvDARTMAE\eV{} and
\PRMDopantCvDARTMAE\eV{}, respectively.  The predefined chemistry block had MAE
\PRMChemistryBlockDARTMAE\eV{} but samples a different target distribution and
should therefore be interpreted separately rather than placed on a monotonic
extrapolation scale.  Across all five regimes, paired absolute-error differences
favored DART over both the descriptor model selected by validation MAE within
each split and the additive-readout SchNet baseline.
Figure~\ref{fig:transfer} places the five transfer regimes and their paired
comparator contrasts on a common visual scale; Table~\ref{tab:benchmark}
reports the complete intervals.

% Auto-generated; do not edit.
\begin{table*}[t]
\caption{Pooled out-of-fold or single-block test MAE (eV). The descriptor
column uses the family selected separately within each regime by validation
MAE (Histogram GB, LightGBM in the resulting selections). $\Delta$
is comparator minus DART absolute error with a paired 95\% bootstrap interval
using the regime-matched resampling unit; positive values favor DART.}
\label{tab:benchmark}
\centering
\begin{tabular}{lccc cc}
\toprule
Regime & DART & SchNet & Descriptor & $\Delta_{\rm SchNet}$ & $\Delta_{\rm desc.}$ \\
\midrule
Random OOF & 0.418 & 0.571 & 0.625 & +0.154 [+0.136, +0.174] & +0.208 [+0.194, +0.221] \\
Pair OOF & 0.450 & 0.628 & 0.674 & +0.178 [+0.157, +0.202] & +0.224 [+0.207, +0.242] \\
Host OOF & 0.938 & 6.703 & 1.273 & +5.765 [+3.006, +9.296] & +0.335 [+0.210, +0.461] \\
Impurity OOF & 0.844 & 1.692 & 1.331 & +0.848 [+0.626, +1.092] & +0.487 [+0.363, +0.619] \\
Chemistry block & 0.386 & 0.508 & 0.642 & +0.122 [+0.069, +0.175] & +0.256 [+0.189, +0.323] \\
\bottomrule
\end{tabular}
\end{table*}


\begin{figure*}[t]
\centering
\includegraphics[width=\textwidth]{fig_transfer.pdf}
\caption{Interpolation and chemical-transfer evaluation.
  (a) DART pooled test MAE across the five evaluation regimes; the dashed line
  marks random out-of-fold MAE.
  (b) Paired differences in absolute error for the additive-readout SchNet
  baseline and the descriptor model selected separately by validation MAE in
  each split, relative to DART.  Intervals are 95\%
  bootstrap intervals using the regime-matched resampling units defined in
  Sec.~\ref{sec:evaluation}; the broken axis isolates the large host-held-out
  SchNet contrast.  Positive values favor DART.}
\label{fig:transfer}
\end{figure*}

Aggregate MAE can hide failures in chemically sparse or reranking-relevant
subsets.  Table~\ref{tab:applicability} therefore reports equal-weight host and
impurity macro errors together with performance on nonpositive formation
energies and the lowest-energy decile.  For example, on pair-held-out
predictions, the true low-decile MAE was
\PRMPairCvDARTLowEnergyMAE\eV{} and its recall by the predicted low decile was
\PRMPairCvDARTLowEnergyRecall\%.  None of these low-energy summaries entered
architecture or hyperparameter selection.

% Auto-generated; do not edit.
\begin{table*}[t]
\caption{Applicability-domain diagnostics for DART. Macro MAEs weight each
host or impurity equally. Nonpositive-target MAE uses
$E_\mathrm{f}\leq0$; low-decile MAE and recall use the true and predicted
lowest 10\% of each pooled test regime. Energy errors are in eV and recall is
in percent. A dash denotes a regime with no eligible nonpositive test
targets.}
\label{tab:applicability}
\centering
\begin{tabular}{lccccc}
\toprule
Regime & Host macro & Impurity macro & Nonpositive-target MAE & Low-decile MAE & Low-decile recall \\
\midrule
Random OOF & 0.489 & 0.440 & 0.396 & 0.440 & 87.1 \\
Pair OOF & 0.533 & 0.468 & 0.450 & 0.480 & 86.0 \\
Host OOF & 0.990 & 0.968 & 1.134 & 1.284 & 69.2 \\
Impurity OOF & 0.887 & 0.826 & 0.973 & 1.041 & 71.6 \\
Chemistry block & 0.384 & 0.387 & -- & 0.305 & 75.0 \\
\bottomrule
\end{tabular}
\end{table*}


Incorporation geometry introduced an additional asymmetry.  Under
pair-held-out prediction, interstitial structures had MAE
\PRMPairInterstitialMAE\eV{}, compared with \PRMPairAdsorbateMAE\eV{} for
adsorbates.  Host-, impurity-, and site-resolved errors were also heterogeneous;
the complete descriptive analysis is reported in Supplemental
Fig.~S4~\cite{supplemental}.

%----------------------------------------------------------------------
\subsection{Internal uncertainty and selective prediction}
\label{sec:uq_results}
%----------------------------------------------------------------------

On the internal ensemble test partition, the five-member ensemble had MAE
\PRMUQTestMAE\eV{}, Gaussian negative log
likelihood \PRMUQNLL, and mean Gaussian CRPS \PRMUQCRPS\eV{}.  The Spearman
correlation between calibrated standard deviation and absolute error was
\PRMUQErrorSpearman.  The risk--coverage curve had AURC \PRMUQAURC\eV{} and
excess AURC \PRMUQExcessAURC\eV{} relative to oracle error ordering.

At 90\% nominal coverage, the split-conformal interval covered
\PRMUQNinetyCoverage\% of test structures with mean width
\PRMUQNinetyWidth\eV{}.  Supplemental Table~SII and
Fig.~S3~\cite{supplemental} report all evaluated nominal levels, the
risk--coverage curve, and the error--uncertainty relation.  These values measure
empirical marginal coverage within the internal ensemble split, not
calibration under a chemically shifted distribution.

%----------------------------------------------------------------------
\subsection{Retrospective out-of-fold reranking within IMP2D}
\label{sec:reranking_results}
%----------------------------------------------------------------------

The host--impurity-pair folds provide one out-of-fold prediction for every
retained structure.  Both the input geometries and reference energies come
from the completed IMP2D DFT workflow, so the following comparisons are
retrospective rather than a substitute for relaxation.  Across pairs
containing both incorporation classes, DART
identified the lower-energy class among the retained, DFT-relaxed IMP2D
candidates with
\PRMPreferenceAccuracy\% accuracy; the predicted and reference class margins
had Spearman correlation \PRMPreferenceMarginSpearman{} and margin MAE
\PRMPreferenceMarginMAE\eV{}.  Selecting one configuration globally gave
\PRMGlobalExactAccuracy\% exact-site accuracy and mean regret
\PRMGlobalRegret\eV{}.  When the incorporation class was fixed, exact and
top-two site accuracies were \PRMWithinTypeExactAccuracy\% and
\PRMWithinTypeTopTwoAccuracy\%, respectively, with mean regret
\PRMWithinTypeRegret\eV{}.
The corresponding descriptive reference accuracies were
\PRMPreferenceReferenceAccuracy\%, \PRMGlobalReferenceAccuracy\%,
\PRMWithinTypeReferenceExactAccuracy\%, and
\PRMWithinTypeReferenceTopTwoAccuracy\%, respectively.  Relative to these
references, DART improved accuracy by
\PRMPreferenceGain{} [\PRMPreferenceGainLow, \PRMPreferenceGainHigh],
\PRMGlobalAccuracyGain{} [\PRMGlobalAccuracyGainLow,
\PRMGlobalAccuracyGainHigh], \PRMWithinTypeExactAccuracyGain{}
[\PRMWithinTypeExactAccuracyGainLow, \PRMWithinTypeExactAccuracyGainHigh], and
\PRMWithinTypeTopTwoAccuracyGain{} [\PRMWithinTypeTopTwoAccuracyGainLow,
\PRMWithinTypeTopTwoAccuracyGainHigh] percentage points, respectively.

Complete decision metrics and clustered intervals are reported in
Supplemental Table~SIII~\cite{supplemental}.  Figure~\ref{fig:reranking}
shows the class-margin agreement, the distribution of reranking regret, and
the clustered decision-accuracy intervals.

\begin{figure*}[t]
\centering
\includegraphics[width=\textwidth]{fig_screening.pdf}
\caption{Retrospective reranking of available DFT-relaxed IMP2D geometries
  using pair-held-out predictions.
  (a) Predicted versus reference preference margins between adsorbate and
  interstitial configurations.
  (b) Empirical cumulative distributions of global and within-class reranking
  regret.
  (c) Host--impurity-pair cluster-bootstrap intervals for incorporation-class,
  exact-site, and top-two decision accuracy.}
\label{fig:reranking}
\end{figure*}

\else

\noindent\textit{Numerical results are unavailable because the audited result
bundle is not present.}

\fi
